\section{Related work}\label{sec:related}

Randomization-based causal inference treats the assignment mechanism as the source of uncertainty and potential outcomes as fixed features of the finite population. This design-based perspective goes back to \citet{fisher_1935}, \citet{neyman_1990}, and the potential-outcomes formulation of \citet{rubin_1974,rubin_1978}; see also \citet{holland_1986}, \citet{imbens_rubin_2015}, \citet{imbens_wooldridge_2009}, \citet{athey_imbens_2017}, and \citet{ding_2024}. Recent econometric treatments emphasize how the choice of randomization affects precision, robustness, and finite-sample behavior, including regression adjustment and covariate balance \citep{freedman_2008,lin_2013,dasgupta_2015,morgan_2012,li_2018,bai_2024}. The present paper works in the same finite-population tradition, with a fixed number of treatment arms \(K\), a fixed number of labeled units \(n\), binary potential outcomes, and a nonzero zero-sum treatment contrast \(c\). Its target is the exact design-estimator minimax risk and its first- and second-order behavior under the randomization distribution.

Finite-population sampling theory provides a parallel design-based language for choosing sampling laws and estimators under adversarial finite populations. Classical unequal-probability sampling and Horvitz--Thompson estimation \citep{horvitz_1952,hansen_1953,sarndal_1992} connect naturally to randomized experiments because treatment assignment can be read as a sampling design over exposure cells. \citet{gabler_1990} and \citet{aronow_lopatto_2026_minimax_unbiased} study minimax ideas in finite-population sampling, with the latter giving a close comparison through finite-population minimax sampling under unbiasedness. The analysis here is aligned with that decision-theoretic finite-population viewpoint and specializes it to multi-arm binary randomization, where both the assignment law and the clipped estimator enter the minimax criterion.

The paper is also connected to optimum experimental design. The classical design literature studies allocation rules through criteria derived from information matrices or risk functionals \citep{smith_1918,kiefer_1959,kiefer_1974,atkinson_2007}. In causal experiments, minimax and optimality criteria have been used to evaluate and construct randomization designs under finite-population or model-assisted objectives. \citet{kallus_2018} develops optimal a priori balance and \citet{kallus_2020_optimality} studies optimal randomization designs; related modern work includes \citet{basse_2023}, \citet{harshaw_2024}, \citet{kandiros_2024,kandiros_2026}, \citet{eckles_2014}, \citet{derezinski_2019}, and \citet{hu_2024}. Relative to this literature, the contribution is an exact finite binary minimax analysis for treatment contrasts, including the orbit reduction and finite-program certificates used to evaluate sharp finite-\(n\) quantities.

Information inequalities and local asymptotic arguments supply the lower-bound tools behind many minimax results. The decision-theoretic tradition begins with \citet{wald_1950}; asymptotic minimax and local experiment methods are developed in \citet{lecam_1986}, \citet{ibragimov_1981}, and related work. Bayesian Cramér--Rao and van Trees inequalities \citep{van_trees_1968,gill_1995}, together with classical information identities \citep{rao_1945}, provide lower bounds for estimation risk. Second-order minimax refinements have a long statistical lineage, including \citet{levit_1981_second_order}, and related shrinkage and nonparametric minimax phenomena appear in \citet{bickel_1981}, \citet{pinsker_1980}, \citet{feldman_1991}, and \citet{johnstone_1992}. The results below use this tradition to identify the first-order minimax constant and to analyze the size of the second-order improvement in the finite-population binary randomization game.

The closest econometric comparisons are recent sharp analyses of two-arm designs. \citet{sudijono_dobriban_tchetgen_2026_sharp} obtains a sharp two-arm binary expansion, and \citet{hull_2026_bounded} studies bounded two-arm potential outcomes. \citet{rosenman_hunter_2026_shrinkage} is also related through shrinkage ideas for experimental estimation. The present paper places those two-arm themes inside a multi-arm contrast framework: it gives exact first-order risk for general fixed contrasts, embeds the two-arm converse where it is sharp, and constructs finite linear-program certificates for rational and real contrasts. A three-arm example then shows, at \(n=3\) and for the single contrast \(c^\dagger=(1,-1/2,-1/2)\) under one fixed allocation, that estimators using the full observed arm label and outcome achieve a strictly smaller worst-case risk than estimators using only a scalar signed score. That comparison isolates what the data reduction discards for this fixed three-arm support, fixed allocation, and sample size.
